\documentclass[spanish,12pt]{report}
\makeatletter
\def\input@path{{../}}
\makeatother
\input{core/preambulo.tex}
\input{core/glosario.tex}
\input{core/acronimos.tex}
\graphicspath{{../}{../../common/}}
\usepackage[inline]{enumitem}
\usepackage{listings}

% JSON is not a built-in language in the listings package.
\lstdefinelanguage{json}{
    basicstyle=\ttfamily\small,
    showstringspaces=false,
    breaklines=true,
    literate=
     *{0}{{{0}}}{1} {1}{{{1}}}{1} {2}{{{2}}}{1} {3}{{{3}}}{1}
      {4}{{{4}}}{1} {5}{{{5}}}{1} {6}{{{6}}}{1} {7}{{{7}}}{1}
      {8}{{{8}}}{1} {9}{{{9}}}{1}
      {:}{{{:}}}{1} {,}{{{,}}}{1}
      {\{}{{{\{}}}{1} {\}}{{{\}}}}{1}
      {[}{{{[}}}{1} {]}{{{]}}}{1},
}

%%%%%%%%%%%%%%%%%%%%%%%%%%%%%%%%%%%%%%%%%%%%%%%%%%%%%%%%%%
\newcommand{\nombreautor}{Lucas Zamora Fochesato}
\newcommand{\mes}{Mes}
\newcommand{\anio}{2026}
\newcommand{\titulo}{Análisis del espacio de instancias (ISA) para la programación semidefinida y generación de nuevas instancias\\{\small\colorbox{yellow}{\textcolor{red}{\textbf{(Título provisional)}}}}}
\newcommand{\nombreprofuno}{Claudio Telha}
\newcommand{\codigo}{ING-IN-XXX/XX}

%%%%%%%%%%Aqui se escribe el resumen%%%%%%%%%%%%%%%%%%%%%

\newcommand{\dedicatoria}{
%En esta sección se deben incluir, de ser necesario, los agradecimientos personales a familiares, colegas, etc. (Máximo 1 página)
 }

% \newcommand{\resumen}{
% En esta sección se debe incluir un resumen de los principales aspectos desarrollados en la tesis, justificación, metodología y principales resultados. (Máximo 1 página). 

% \vspace{5cm}

% {\color{blue}

% El siguiente documento no puede ser la copia del escrito del o los art\'{\i}culos publicados (si corresponde).
% Estos deben ser incluidos como Ap\'endices al final de la tesis.

% \vspace{1cm}

% Los m\'argenes son 2.5 cm superior e inferior, y 3 cm derecha e izquierda, en todo el documento.

% \vspace{1cm}

% Se debe mantener el formato tipo de letra Times New Roman, tama\~no 12 e interlineado 1.5 l\'{\i}neas.
% }
% }

\newcommand{\resumen}{

}


\newcommand{\agradecimientos}{
En esta sección, se deben incluir los agradecimientos en primer lugar a las fuentes de financiamiento de las tesis, si las hubiera y posteriormente a los académicos involucrados en el desarrollo de la investigación. (Máximo 1 página). 
}
\renewcommand{\baselinestretch}{1.5}

%%%%%%%%%%%%%%%%%%%%%%%%%%%%%%%%%%%%%%%%%%%%%%%%%%%%%%%%%%

% Por favor coloque todas las definiciones de simbolos a continuacion
% y no en los archivos de capitulos. Esto evitara la  existencia de
% multiples definiciones para las mismas palabras.

\newcommand{\mydef}
	{\stackrel{\mathrm{def}}{=}}
\newcommand{\e}
	{\hbox{\large{e}}}
\newcommand{\mi}
	{\hbox{\large{i}}}
\newcommand{\RR}{\mathbb{R}}


%%%%%%%%%%%%%%%%%%%%%%%%%%%%%%%%%%%%%%%%%%%%%%%%%%%%%%%%%%

\begin{document}
\label{start}

% Genera las paginas de titulo, copyright, resumen,
% agradecimientos, dedicatoria, indices,...

% No es necesario modificar

\input{core/primeras_paginas.tex}

%%%%%%%%%%%%%%%%Ingreso de capitulos%%%%%%%%%%%%%%%%%%%%%%

% espaciado similar a Word 1.5
\setstretch{1.4}

\input{chapters/chapter01.tex} % Intro
\input{chapters/chapter02.tex} % Rev. biblio
\input{chapters/chapter03.tex} % Metodología
\input{chapters/chapter04.tex} % Resultados
\input{chapters/chapter05.tex} % Conclusiones

\cleardoublepage

\printglossary[type=\acronymtype, title=Nomenclatura, toctitle=Nomenclatura]

\printglossary[title=Glosario, toctitle=Glosario]


% Ingenieros prefieren la bibliografia antes que los anexos

\cleardoublepage

% {\color{blue} En esta sección de debe incluir todas las referencias utilizadas en la tesis. Se debe evitar el uso de páginas web o fuentes no formales de información. (Máximo 15 páginas)

% \vspace{0.5cm}

% Sobre el formato de citas, se debe usar el formato APA, esto es: ``Según datos de Fernández (1999), un 110\% de los conductores no alcanza a frenar... En otro estudio del mismo autor (Fernández, 2000) se establece... Otras investigaciones (Fernández y Vergara, 2000) concluyen además... El trabajo de Fernández et al. (2001) confirma que...''

% \vspace{0.5cm}

% Para mayor información consultar: \url{https://uandes.libguides.com/citasyreferencias} 

\printbibliography[heading=bibintoc,title={Referencias}]


\cleardoublepage

\appendix % seguido de archivos de apendices

% Escribe 'Anexos' en el Indice General
\addcontentsline{toc}{chapter}{Anexos}


\input{attachments/anexo_a.tex}    % archivo de apendice

% \input{ruta_al_Archivo/(apendice_b}    % otro archivo de apendice

\label{end}
\end{document}

%%%%%%%%%%%%%%%%%%%%%%%%%%%%%%%%%%%%%%%%%%%%%%%%%%%%%%%%%%%%
